% !TEX program = xelatex
% ============================================================
%  كتاب الوصفات — Recipe Book Template
%  قالب صفحة وصفة: اسم الوصفة، معلومات (وقت التحضير|الطبخ|الحصص|الصعوبة)،
%  مقاديمرتبة في عمودين، خطوات مرقّمة، ملاحظات، وقيم غذائية.
%  Compile with: xelatex main.tex (run twice)
% ============================================================
%  Features:
%    - Recipe name header with decorative rule
%    - Meta row: prep time | cook time | servings | difficulty
%    - Two-column ingredients list
%    - Numbered cooking steps
%    - Notes box (tcolorbox) and nutrition placeholder
%    - Reusable per recipe (sample recipe included)
% ============================================================

\documentclass[11pt,a4paper]{article}

% ---- Geometry ----
\usepackage[a4paper,margin=2cm,top=2cm,bottom=2cm,headheight=16pt]{geometry}

% ---- Core packages (order matters for bidi) ----
\usepackage{graphicx}
\usepackage{array}
\usepackage{booktabs}
\usepackage{tabularx}
\usepackage{multirow}
\usepackage{enumitem}
\usepackage{multicol}
\usepackage{float}
\usepackage{titlesec}
\usepackage{fancyhdr}
\usepackage{setspace}
\usepackage{xcolor}
\usepackage{tcolorbox}
\tcbuselibrary{skins,breakable}
\usepackage{tikz}

% ---- RTL / Arabic language support (must load before hyperref) ----
\usepackage{bidi}
\usepackage[no-math]{fontspec}
\usepackage[quiet]{polyglossia}

% ---- Fonts ----
\setmainfont[Scale=1]{NotoKufiArabic-Regular.ttf}
\newfontfamily\arabicfont[Script=Arabic,Scale=0.95]{NotoKufiArabic-Regular.ttf}
\newfontfamily\englishfont[Scale=0.95]{TeX Gyre Termes}

% ---- Languages ----
\setdefaultlanguage[calendar=gregorian,hijricorrection=1]{arabic}
\setotherlanguage[variant=british]{english}

% ---- Hyperref (load after polyglossia/bidi) ----
\usepackage[breaklinks,hidelinks]{hyperref}
\hypersetup{
  pdftitle={كتاب الوصفات},
  pdfauthor={الاسم}
}

% ---- Colors ----
\definecolor{recipeaccent}{HTML}{8B3A2A}
\definecolor{recipebg}{HTML}{FBF3EE}
\definecolor{recipeink}{HTML}{2B2B2B}

% ---- Line spacing ----
\setstretch{1.3}

% ---- Section formatting ----
\titleformat{\section}{\Large\bfseries\color{recipeaccent}}{\thesection}{0.7em}{}
\titleformat{\subsection}{\large\bfseries\color{recipeaccent}}{\thesubsection}{0.6em}{}

% ---- Header / footer ----
\pagestyle{fancy}
\fancyhf{}
\fancyhead[R]{\small كتاب الوصفات}
\fancyhead[L]{\small الاسم}
\fancyfoot[C]{\small \thepage}
\renewcommand{\headrulewidth}{0.4pt}

% ---- Recipe meta cell ----
\newcommand{\metacell}[2]{%
  \begin{center}
    {\small\color{recipeaccent}\textbf{#1}}\\[2pt]
    {\normalsize #2}
  \end{center}}

% ---- Notes box ----
\newtcolorbox{notesbox}{
  colback=recipebg,colframe=recipeaccent,boxrule=0.6pt,arc=2pt,
  fonttitle=\bfseries,title=ملاحظات,
  left=8pt,right=8pt,top=6pt,bottom=6pt,breakable
}

\begin{document}
\color{recipeink}

% ============================================================
%  Recipe: وصفة تجريبية (Sample Recipe)
%  TODO: انسخ كتلة الوصفة هذه لكل وصفة جديدة وعدّل المحتوى.
% ============================================================

\section*{وصفة تجريبية (Sample Recipe)}
% TODO: استبدل باسم الوصفة الحقيقية
{\large\itshape\color{recipeaccent}
وصفة عربية تقليدية — مثال توضيحي\par}

\vspace{0.2cm}
\begin{tikzpicture}
  \draw[recipeaccent,line width=1pt] (0,0) -- (\linewidth-1cm,0);
\end{tikzpicture}
\vspace{0.4cm}

% ---- Meta row ----
\noindent
\begin{tabularx}{\linewidth}{|X|X|X|X|}
\hline
\metacell{وقت التحضير}{٢٠ دقيقة} &
\metacell{وقت الطبخ}{٤٥ دقيقة} &
\metacell{الحصص}{٤ أشخاص} &
\metacell{الصعوبة}{متوسطة} \\
\hline
\end{tabularx}

\vspace{0.6cm}

% ---- Ingredients ----
\subsection*{المقادير (Ingredients)}
% TODO: املأ قائمة المقادير
\begin{multicols}{2}
\begin{enumerate}[label=\textarabic{\arabic*.},leftmargin=1.4em,itemsep=2pt]
  \item كوبان من الدقيق
  \item كوب من السكر
  \item نصف كوب من زيت الزيتون
  \item ثلاث بيضات
  \item ملعقة صغيرة من البيكنج باودر
  \item رشة ملح
  \item ملعقة صغيرة من الفانيليا
  \item كوب من الحليب
\end{enumerate}
\end{multicols}

\vspace{0.4cm}

% ---- Steps ----
\subsection*{طريقة التحضير (Instructions)}
% TODO: املأ خطوات التحضير
\begin{enumerate}[label=\textarabic{\arabic*.},leftmargin=1.6em,itemsep=4pt]
  \item يُسخّن الفرن إلى ١٨٠ درجة مئوية وتُدهن صينية الخبز.
  \item يُنخل الدقيق مع البيكنج باودر والملح في وعاء كبير.
  \item يُخفق البيض مع السكر والفانيليا حتى يصبح المزيج كريميًا.
  \item يُضاف زيت الزيتون والحليب تدريجيًا مع الاستمرار في الخفق.
  \item يُضاف خليط الدقيق تدريجيًا ويُحرّك حتى تتجانس المكونات.
  \item يُصب الخليط في الصينية وتُخبز لمدة ٤٠–٤٥ دقيقة.
  \item تُترك لتبرد قبل التقديم.
\end{enumerate}

\vspace{0.5cm}

% ---- Notes box ----
\begin{notesbox}
% TODO: أضف ملاحظاتك هنا
\begin{itemize}[leftmargin=1.2em,itemsep=2pt]
  \item يمكن إضافة قشر الليمون المبشور لنكهة إضافية.
  \item يُحفظ الكيك في وعاء محكم لمدة تصل إلى ثلاثة أيام.
\end{itemize}
\end{notesbox}

\vspace{0.5cm}

% ---- Nutrition placeholder ----
\subsection*{القيمة الغذائية (Nutrition)}
% TODO: املأ القيم الغذائية لكل حصة
\noindent
\begin{tabularx}{\linewidth}{|X|X|X|X|}
\hline
\textbf{السعرات} & \textbf{الدهون} & \textbf{الكربوهيدرات} & \textbf{البروتين} \\
\hline
٢٥٠ سعرة & ١٢ جم & ٣٠ جم & ٦ جم \\
\hline
\end{tabularx}

\end{document}
